\section{An ordinary alternating-core shape restriction}
\label{sec:ordinary-alternating-shapes}

Spin representations have been removed by HD-SPIN-02.  Ordinary
alternating representations require a different local estimate:
restriction to the six-letter subgroup can contain trivial summands.
We retain their contribution explicitly.

\begin{lemma}[HD-ALT-01a: the six-letter fixed-space estimate]
\label{HD-ALT-01a}
Let $N=\mathfrak A_n$, $n\ge8$, act faithfully and irreducibly
on $W$, with character $\chi$ and degree $a$.  Let $c$ be a
three-cycle and put $u=\chi(c)/a$.  Every nonscalar element $g$
of the linear normalizer of $N$, of finite order up to a scalar,
satisfies
\[
                \nu(g)\ge\frac{2a(1-u)}{15}.              \tag{ALT.1}
\]
Here $u$ is real and belongs to $[-1/2,1)$.
\end{lemma}

\begin{proof}
The complete ordinary character table of $\mathfrak A_6$ gives
\[
       \dim\ker(h-\lambda)\le\frac35\dim U              \tag{ALT.2}
\]
for every nontrivial irreducible $U$, every $h\ne1$ and every
$\lambda$.  We record the entire finite scope of this input.
The six nontrivial center-positive rows of the retained table
$2.A6$, which factor through $A6$, have degrees
$5,5,8,8,9,10$.  The minima of $\nu(h)/\dim U$ across all
noncentral lift classes are respectively
\[
                   2/5,\ 2/5,\ 1/2,\ 1/2,\ 4/9,\ 2/5.
                                                               \tag{ALT.3}
\]
The exact Fourier formula (SPINQ.3) checks all $66$ spectra,
including every possible eigenvalue.  The full spectra,
nonnegative integral multiplicities and degree sums are retained
in \texttt{research/alternating\_local\_certificate.json}.
The character table and identification with the named group
are published inputs \cite{SPOR:CTblLib}; the finite calculation
proves the claimed numerical inequality on those complete inputs.

As in HD-SPIN-01a, the normalizer modulo scalars embeds in
$\operatorname{Aut}(\mathfrak A_n)=\mathfrak S_n$.
For a nonscalar $g$ choose a three-cycle $c'$ not centralized by
its induced permutation, and put $h=[g,c']$.  This is nonidentity
and supported on at most six letters.  Embed its support in a
six-letter set and let $A6$ be the corresponding natural subgroup.
All such subgroups are conjugate in $\mathfrak A_n$ for $n\ge8$:
an odd conjugating permutation can be corrected by a transposition
on two letters outside the target six-set.  Thus the dimension
$d_0=\dim W^{A6}$ is independent of this choice.

On restriction, write $W=\mathbf1^{d_0}\oplus W'$, with no
trivial constituent in $W'$.  For every fixed eigenvalue,
(ALT.2) adds across its irreducible constituents.  The trivial
part contributes at most $d_0$, whence
\[
                    \nu(h)\ge\frac25(a-d_0).
\]
The scalar-adjusted commutator inequality gives
$\nu(h)\le2\nu(g)$, and therefore
$\nu(g)\ge(a-d_0)/5$.

The $A6$-invariants are contained in the fixed space of any
three-cycle in that subgroup.  A three-cycle is conjugate to
its inverse in $\mathfrak A_n$ (use a transposition on its
support and another on two fixed letters).  Consequently
$\chi(c)$ is real, and its two nontrivial cube-root eigenvalues
have equal multiplicity.  Its fixed dimension is
$(a+2\chi(c))/3$, so
\[
              d_0\le\frac{a+2\chi(c)}3.
\]
Combining the inequalities proves (ALT.1).  Nonnegativity of the
fixed dimension gives $u\ge-1/2$, and faithfulness gives $u<1$.
\end{proof}

\begin{lemma}[HD-ALT-01b: the Young-diagram bound]
\label{HD-ALT-01b}
Let $n\ge15$ and let $\chi$ be an irreducible ordinary character
of $\mathfrak A_n$ occurring in the restriction of the Specht
character indexed by $\lambda\vdash n$.  If
\[
               \max(\lambda_1,\lambda'_1)\le\lceil n/2\rceil,
                                                               \tag{ALT.4}
\]
then its three-cycle ratio satisfies
\[
             u\le\frac{n-5}{4(n-2)}
               =\frac14-\frac3{4(n-2)}.                  \tag{ALT.5}
\]
For even $n$ the sharper numerator is $n-7$.
\end{lemma}

\begin{proof}
We use the published Jucys--Murphy content-eigenvalue theorem
\cite[Theorem~4.2; author version p.~9]{HD:CMN2022}.
For $J_i=\sum_{j<i}(j\ i)$, direct multiplication in
$\mathbf C[\mathfrak S_n]$ gives
\[
       \sum_iJ_i^2=\binom n2\,1+
                  \sum_{c\text{ a three-cycle}}c.
\]
Indeed the equal-transposition products give $i-1$ identities,
and the ordered unequal pairs give each oriented three-cycle
exactly once, indexed by its largest letter.  The content theorem
therefore gives the following short specialization:
\[
 u=\frac{3\left(\sum_{(i,j)\in\lambda}(j-i)^2-\binom n2\right)}
           {n(n-1)(n-2)}.                                \tag{ALT.6}
\]
This also applies if the Specht restriction splits.  The two
constituents have equal degrees and are interchanged by odd
conjugation.  The three-cycle class does not split in
$\mathfrak A_n$, since a transposition on two fixed letters
centralizes the cycle.  Their character values on it are equal,
so their normalized ratio is the Specht ratio in (ALT.6).
The restriction rule itself is the standard one explicitly
recorded in \cite[Section~2A, p.~604]{HD:BNOT2015}.

Put $k=\lceil n/2\rceil$ and use Frobenius coordinates
$(a_1,\ldots,a_d\mid b_1,\ldots,b_d)$ of $\lambda$.
The diagonal hooks partition its boxes, with contents
$-b_i,\ldots,-1,0,1,\ldots,a_i$ on the $i$th hook.
Condition (ALT.4) gives $0\le a_i,b_i\le k-1$ and
$\sum_i(a_i+b_i)=n-d\le n-1$.  For $0\le s\le k-1$,
\[
 \sum_{j=1}^s j^2=\frac{s(s+1)(2s+1)}6
                  \le\frac{s\,k(2k-1)}6.
\]
Hence the content-square sum is at most
$k(2k-1)(n-1)/6$.  Substitution in (ALT.6) gives
\[
                u\le\frac{k(2k-1)-3n}{2n(n-2)}.
\]
For even $n$ this equals $(n-7)/(4(n-2))$, and for odd $n$
it equals $(n-5)/(4(n-2))$.  This proves both assertions.
\end{proof}

\begin{theorem}[HD-ALT-01c: exclusion of the indicated ordinary shapes]
\label{HD-ALT-01c}
Suppose a primitive inducing block of a dominant nonsplit Weil
sixfold source is essentially self-dual and has an irreducibly
acting normal core $E=\mathfrak A_n^t$, with $n\ge15$.
Suppose every irreducible tensor factor has a Specht label
satisfying (ALT.4).  Then these data cannot occur.

In particular, for the transitive ordinary alternating core of
D-BI-19, a remaining essentially self-dual source in alternating
degree $n\ge15$ must have
$\max(\lambda_1,\lambda'_1)>\lceil n/2\rceil$.
All self-dual constituents of self-conjugate labels are excluded
in this range.  Nonselfdual constituents are not covered.
\end{theorem}

\begin{proof}
Let $e$ be the total inducing degree.  Lemmas HD-ALT-01a/b give,
for each factor-preserving nonscalar normalizer element,
\[
        \nu(g)\ge\frac{n-1}{10(n-2)}e.                   \tag{ALT.7}
\]
The same conclusion holds for nontrivial factor permutations.
Indeed the intrinsic factor-algebra argument and the weighted
cyclic-word estimate (SPOR.6) give $\nu(g)\ge e/4$ in that
case, and $(n-1)/(10(n-2))<1/4$.  Genuine scalar inertia moves
all $e$ dimensions.  This argument permits every compatible
outer automorphism and all tensor permutations.

The actual intermediate-curve and nontrivial-similitude argument
of D-BI-09 gives at least twelve genuine inertias, of which at
least two move at least $\lceil e/2\rceil$ dimensions.  Its
hypotheses hold here: the source has nonsplit actual polarization,
exact imaginary quadratic original character field, and a
primitive essentially self-dual inducing block; its inducing
character field may be larger.  The other ten inertias obey
(ALT.7).  The continuous consequence of the exact integer
budget is
\[
 e+\frac{n-1}{n-2}e\le2e+6,
                    \qquad e\le6(n-2).                  \tag{ALT.8}
\]
The factor degrees are at least $n-1$ in this range by
\cite[Lemmas~3.1--3.2, p.~607]{HD:BNOT2015}; in particular
$e\ge9$, as required in D-BI-09.

Choose any one factor character and its label $\lambda$.
The degree of the corresponding $\mathfrak S_n$ character
is either its factor degree or twice that degree, so by (ALT.8)
it is at most $12(n-2)$.  For $n\ge15$,
\[
       12(n-2)<\frac{2n^3}{27}.
\]
To see this uniformly, the difference after multiplication by
$27$ is $2n^3-324n+648$; at $n=15$ it is $2538$, and its forward
difference is $6n^2+6n-322>0$ for $n\ge15$.
The published precise shape consequence of Rasala's classification
\cite[proof of Theorem~9.7, p.~312]{HD:GKT2019}, stated there
for every $\mathfrak S_n$ character of degree
$2\le d<2n^3/27$, therefore forces $\lambda$, up to transpose,
to be one of
\[
             (n-1,1),\qquad(n-2,2),\qquad(n-2,1,1).
\]
Each of these shapes has a row or column of length at least
$n-2>\lceil n/2\rceil$, contrary to (ALT.4).

Finally, if a partition is self-conjugate then
$2\lambda_1-1\le n$, so (ALT.4) holds automatically.
The preceding conclusion thus excludes its self-dual constituents.
For the transitive core of D-BI-19, factor labels agree up to the
outer action and transpose, which preserves (ALT.4); hence the
stated necessary shape restriction applies to that gate.
\end{proof}

\paragraph{Scope.}
The new local calculation and shape inequality do not prove
an all-alternating exclusion.  Long first-row or first-column
shapes and nonselfdual split characters remain.  The already
excluded standard, exterior-square and unordered-pair families
retain their stronger independent all-degree statements.
